% FailWhere — Dataset documentation supplement (DMLR submission)
% Follows the NeurIPS Datasets & Benchmarks documentation guidance.
\documentclass[11pt]{article}
\usepackage[margin=1in]{geometry}
\usepackage{url}
\usepackage{enumitem}

\title{FailWhere: Dataset Documentation and Intended Uses\\
\large Supplementary material}
\author{Subodh Wandile}
\date{}

\begin{document}
\maketitle

\section*{1. Overview}
\textsc{FailWhere} is a benchmark for testing whether a vision-language
model (VLM) can verbally identify which spatial region of an image caused
its own incorrect answer. The release contains: (a)~a synthetic
grid-image benchmark, (b)~a real-image extension built on the Visual
Spatial Reasoning (VSR) dataset, (c)~counterfactual-masking ground truth
(CF-GT) for 142 model failures, and (d)~the complete model response logs
underlying every number in the paper.

\section*{2. Motivation}
The benchmark was created to isolate \emph{verbal failure
self-localization}: given an answer the model got wrong, can it name the
responsible region using only its verbal output? Prior attribution work
either requires model internals or only evaluates correct answers. The
benchmark is intended for evaluating introspection, faithfulness, and
spatial reasoning of VLMs. It is not intended for model training, and it
is not a measure of general VQA ability.

\section*{3. Composition}

\subsection*{3.1 Synthetic split (\texttt{code/data/qa\_samples\_500.json})}
500 samples. Each image is a 300$\times$300\,px synthetic grid divided
into nine labeled cells (top-left through bottom-right); each cell holds
one shape$\times$color object. Two question types:
\emph{adjacent\_distractor} and \emph{relative\_position}. Ground truth is
the \emph{answer-cell} (cell containing the correct answer), known by
construction.

\subsection*{3.2 VSR split (real images)}
The first 500 items of \texttt{cambridgeltl/vsr\_random}, split
\texttt{test} (Liu et al.\ 2023): COCO train2017 photographs paired with
binary True/False spatial captions. Model outputs and region tags are
released; the underlying images are redistributed only by reference to
the VSR/COCO sources (Section~6).

\subsection*{3.3 Counterfactual-masking ground truth (CF-GT)}
For each of 142 model failures on the VSR split, each of the nine grid
cells was masked with solid gray (128,\,128,\,128) and the model re-run.
A cell whose masking flips the answer to correct is CF-GT; outcomes are
labeled \texttt{clean} (exactly one cell, $n{=}22$),
\texttt{multiple} ($n{=}47$), or \texttt{none\_found} ($n{=}73$).

\subsection*{3.4 Model response logs}
Six result JSONs under \texttt{code/experiments/} record, per item, the
question, model answer, all three self-localization prompt responses
(neutral, stateless, multi-turn), the parsed \textsc{region:} tag, CF-GT,
and per-cell masking answers. A 141-item Sonnet control file
(\texttt{real\_vsr\_v3\_t0\_sonnet\_explain.json}) records a second
model's attributions of Haiku's failures.

\section*{4. Collection and annotation}
No human annotators were used. Synthetic items were generated
programmatically (\texttt{code/generate\_synthetic.py}); ground truth is
therefore exact by construction. VSR items are unchanged from the source
dataset. CF-GT labels were produced mechanically by the masking
procedure, not by human judgment. Model generations were produced via the
Anthropic API: \texttt{claude-haiku-4-5-20251022} /
\texttt{global.anthropic.claude-haiku-4-5-20251001-v1:0} (temperature 0.0)
and \texttt{claude-sonnet-5-5} /
\texttt{global.anthropic.claude-sonnet-5-5} (default temperature; the
Bedrock endpoint rejects the parameter for this model). The COCO
evidence-cell analysis joins VSR captions to COCO
\texttt{instances\_train2017} bounding boxes via a documented noun-matching
heuristic (93.4\% join rate; a known limitation).

\section*{5. Recommended uses and misuses}
\textbf{Recommended:} evaluating verbal self-attribution, faithfulness,
and region-grounding of VLMs; counterfactual evaluation methodology
research. \textbf{Not recommended:} training; measuring general VQA
accuracy; drawing conclusions about models' internal (non-verbal)
representations, which this benchmark does not access.

\section*{6. Distribution, rights, and licensing}
Repository: \url{https://github.com/subodhwandile/vlm-failure-bench},
released under the MIT license. The VSR dataset is a public research
dataset built on COCO (which carries its own per-image licenses); this
release references those images by the source URLs rather than
redistributing them, and users must comply with the COCO terms of use.
Model-generated outputs were produced through a paid API and are released
by the author. The author bears all responsibility in case of any
violation of rights, and confirms the licenses of all released artifacts.

\section*{7. Maintenance plan}
The benchmark is hosted on GitHub and maintained by the author via the
repository issue tracker. Versioned result files are committed with the
code; any future data corrections will be issued as tagged releases and
documented in the repository changelog. No fixed update cadence is
planned; the dataset is a frozen experimental artifact corresponding to
the paper's reported results.

\end{document}
